\section{Introducción}\label{sec:intro}

Las organizaciones reparten a sus usuarios entre canales de calidad y coste distintos: un centro de atención especializado frente a uno generalista, un hospital frente a atención primaria, un gestor personal frente a una aplicación. Antes de redistribuir, quieren saber cuánto se perdería al mover a un grupo hacia el canal básico. Un experimento aleatorizado lo respondería, pero suele ser caro, lento o inviable; lo disponible son datos observacionales de cada canal. El enfoque estándar repondera a las unidades que hoy atiende el canal básico para que se parezcan al grupo objetivo, con pesos estimados por un clasificador. Formalmente, es un cambio de medida.

Este artículo trata de un rasgo de los datos de canal que el enfoque estándar ignora: \emph{cada canal registra las covariables a su manera}. Un gestor personal apunta un historial comercial que la aplicación nunca guarda; un hospital hace pruebas que atención primaria no hace; un canal registra lo que el cliente declara y otro lo mide. Mostramos que esto rompe la reponderación de dos formas, una visible y otra invisible.
\begin{itemize}
  \item \textbf{Una variable que solo registra un canal} rompe el soporte por construcción, por comparables que sean las unidades: la masa sin soporte es al menos la fracción del grupo objetivo con dato, y el clasificador ideal lo delata al estar seguro ($p=1$) en esos perfiles (Proposición~\ref{prop:rompe}); en una referencia experimental (NSW frente a CPS), una variable de ingresos censurada rompe el soporte en su extremo y los clasificadores estimados no lo detectan (sección~\ref{sec:lalonde}). Descartar la variable restaura la identificación cuando no predice el resultado en el canal básico dadas las demás covariables (Proposición~\ref{prop:descartar}); si lo predice, ni conservarla ni descartarla identifica el estimando.
  \item \textbf{Una variable medida con un instrumento distinto en cada canal} puede conservar el soporte. El estimador converge entonces a un valor desplazado, que derivamos exactamente, y que ningún análisis de los datos observados distingue de la verdad (Teorema~\ref{teo:medicion}). Con datos públicos en los que existen los dos instrumentos para las mismas personas (índice de masa corporal declarado y medido en NHANES), el desplazamiento derivado coincide de forma aproximada con el sesgo observado en 7 de 8 combinaciones de diseño y resultado (la excepción, la tensión con LightGBM, se explica por el modelo de resultado), y los diagnósticos habituales de solapamiento mejoran, respecto de la medición correcta, en lugar de delatarlo (sección~\ref{sec:nhanes}).
\end{itemize}

Para enunciar estos resultados hace falta la maquinaria de la reponderación en lenguaje de teoría de la medida (sección~\ref{sec:identificacion}). Casi toda es conocida y no la reclamamos: el estimando es la media contrafactual que entra en un efecto medio sobre los tratados, los pesos de odds $p/(1-p)$ son los de \citet{hirano2003}, estimar cocientes de densidades con clasificadores es habitual \citep{qin1998,bickel2009,sugiyama2012}, y la teoría doblemente robusta con \emph{cross-fitting} es la de \citet{chernozhukov2018}. Dos lecturas son nuevas, y las secciones~\ref{sec:canal} y~\ref{sec:instrumentos} se apoyan en ellas. La parte sin soporte de la ley objetivo es exactamente su restricción a $\{p=1\}$, que es la construcción de von Neumann de la descomposición de Lebesgue \citep[teorema 6.10]{rudin1987} leída en el clasificador; y sin soporte, el estimador doblemente robusto rellena esa parte en silencio extrapolando el modelo de resultado (Teorema~\ref{teo:lebesgue} y Corolario~\ref{cor:drsinsoporte}).

Damos además herramientas para la práctica. Cuando el grupo objetivo es un corte en una puntuación conocida, los pesos son la propensión de canal ordinaria con el corte impuesto (Proposición~\ref{prop:corte}), como en el análisis por subgrupos \citep{rassen2012}; estimar esa propensión con todas las unidades del canal premium redujo la dispersión hasta un 59\,\% en nuestra simulación cuando el solapamiento era bueno, pero no cuando era malo. Damos el sesgo exacto del recorte por tope de los pesos verdaderos con normalización (Proposición~\ref{prop:recorte}), que complementa a \citet{ma2020}, y adaptamos las cotas del modelo de sensibilidad marginal \citep{tan2006,dorn2023} a este estimando, con una calibración mediante el modelo de canal (sección~\ref{sec:sensibilidad}). Además de los dos análisis con datos públicos, una simulación contrasta los demás resultados: en cada escenario calculamos el límite de cada estimador, y registramos otras predicciones, antes de ejecutarlo, e informamos de las predicciones que fallaron (secciones~\ref{sec:diseno} y~\ref{sec:resultados}). El recorte con pesos estimados y el comportamiento del bootstrap son el objeto de una nota complementaria en preparación; aquí solo resumimos los resultados.

La sección~\ref{sec:planteamiento} plantea el problema y la sección~\ref{sec:identificacion}, la maquinaria de la reponderación. Las secciones~\ref{sec:canal} y~\ref{sec:instrumentos} contienen los dos fallos, cada uno seguido de datos públicos: el visible con la referencia NSW/CPS (sección~\ref{sec:lalonde}) y el invisible con NHANES (sección~\ref{sec:nhanes}). Las secciones~\ref{sec:corte}--\ref{sec:sensibilidad} reúnen herramientas para la práctica que no dependen del problema de medición: pesos con corte conocido, recorte e inferencia, y cotas de sensibilidad. Las secciones~\ref{sec:diseno} y~\ref{sec:resultados} presentan la simulación, y la sección~\ref{sec:discusion} concluye.

Usamos el lenguaje de la teoría de la medida (derivadas de Radon--Nikodym) en lugar de cocientes de densidades. No es una aportación, sino la formulación estándar; permite tratar covariables mixtas (continuas, categóricas y con valores ausentes) sin suponer que existen densidades.

\paragraph{Trabajo relacionado.} El mismo estimando, los pesos de odds y la misma función de influencia doblemente robusta aparecen en la literatura sobre generalizar y transportar resultados de ensayos a una población objetivo \citep{westreich2017,dahabreh2019}. \citet{heckman1997} encontraron empíricamente que el emparejamiento elimina gran parte del sesgo cuando el grupo de comparación procede del mismo mercado y responde el mismo cuestionario que los participantes; el Teorema~\ref{teo:medicion} formaliza esa lección como un resultado de indistinguibilidad. \citet{cook2008} y \citet{steiner2011} documentan también la importancia de medir las covariables igual y con fiabilidad, y que el error de medición diferencial no se corrige sin información externa es conocido en la literatura de error de medición; \citet{imai2010} estudian la identificación no paramétrica y el análisis de sensibilidad con error de medición diferencial del tratamiento. Lo más cercano al Teorema~\ref{teo:medicion} es la literatura sobre error de medición de covariables en métodos de puntuación de propensión. \citet{hong2017} tratan una covariable medida de forma distinta en los dos grupos de tratamiento y la corrigen con modelos bayesianos cuyas distribuciones a priori sobre el error actúan como parámetros de sensibilidad; \citet{webbvargas2017} imputan la covariable verdadera a partir de una muestra externa de calibración en la que se miden las dos versiones; y \citet{rudolph2016} adaptan análisis de sensibilidad para la confusión no observada al error de medición de covariables. Estos métodos corrigen el sesgo a partir de un modelo del error o de datos externos. Nuestros resultados son complementarios: caracterizamos el límite de los estimadores ponderado y doblemente robusto con instrumentos propios de cada canal, mostramos que sin esa información el sesgo no está identificado, que es por lo que esas correcciones necesitan distribuciones a priori o datos de validación, tratamos el caso límite de una variable que solo registra un canal, y mostramos que los diagnósticos habituales de solapamiento pueden mejorar a medida que crece el sesgo. Sobre covariables parcialmente observadas, el enfoque por patrón de ausencia \citep{dagostino2000,blake2020} exige que la variable no confunda allí donde falta; nuestras variables de canal son el caso en que el patrón lo fija el tratamiento, y la condición que usamos es un criterio de selección de covariables por el lado del resultado \citep{deluna2011,vanderweele2011}. Sobre diagnósticos, \citet{westreich2011} muestran que un AUC alto del modelo de propensión no es necesario ni suficiente para controlar la confusión; nuestra simulación añade que tampoco mide el solapamiento. Para diagnosticar fallos de positividad, véanse \citet{petersen2012} y \citet{damour2021}. Para el solapamiento limitado y el recorte, véanse \citet{crump2009}, \citet{ma2020,ma2023}, \citet{liu2023} y \citet{manamkoong2025}, cuya extrapolación de la parte sin solapamiento bajo un supuesto de suavidad complementa el Corolario~\ref{cor:drsinsoporte}; los comentamos donde aparecen (sección~\ref{sec:inferencia}).
